\begin{claim}{D.2.14}{Proved}{The odd-primary Pr\"ufer ideal is the full primary ideal}
Let
\[
 D=\Sigma^\infty\mathbb{C}P^\infty,
 \qquad J_p=\operatorname{loc}(D\wedge S/p),
\]
where the suspension spectrum is reduced. The following statements hold.
\begin{enumerate}
\item If an $E_1$-ring $R$ satisfies $F(D,R)\simeq0$ and the image
  $\eta_R\in\pi_1R$ of the stable Hopf class is zero, then $R\simeq0$.
\item For every prime $p$, the inclusion
  $J_p\subseteq\operatorname{loc}(S/p)$ is proper if and only if there
  is a nonzero derived $p$-complete $E_\infty$-ring $R$ such that
  $F(D,R)\simeq0$.
\item For every odd prime $p$,
  \[
   \operatorname{loc}\bigl(\Sigma^\infty B C_{p^\infty}\bigr)
      =J_p=\operatorname{loc}(S/p).
   \tag{D.2.14a}
  \]
  In particular this ideal contains the finite Moore spectrum $S/p$.
\end{enumerate}
The first statement uses an actual trivialization of the coefficient
line bundle. It does not infer Eilenberg--Moore convergence from the
finiteness of $S^1$.
\end{claim}

\begin{proof}
\emph{The Euler sequences and the Thom fiber.}
Put $G=S^1$, and let $V=\mathbb C$ be its standard complex
representation. The equivariant Euler cofiber sequence is
\[
 \Sigma^\infty G_+\longrightarrow S\longrightarrow S^V.
 \tag{D.2.14b}
\]
Taking homotopy orbits gives
\[
 S\longrightarrow\Sigma^\infty BG_+
       \longrightarrow\operatorname{Th}(L),
\]
where $L$ is the universal complex line over $BG=\mathbb CP^\infty$.
The first map comes from $EG\times_GG=EG\to BG$; since $EG$ is
contractible, it is the inclusion of a basepoint up to homotopy.
Consequently $\operatorname{Th}(L)\simeq D$.

Tensoring \textup{(D.2.14b)} with $S^{-V}$ before taking homotopy
orbits instead gives
\[
 S^{-2}\xrightarrow{j}\operatorname{Th}(-L)
   \longrightarrow\Sigma^\infty BG_+
   \xrightarrow{\delta}S^{-1}.
 \tag{D.2.14c}
\]
Indeed $G_+\wedge S^{-V}$ is induced from the underlying sphere
$S^{-2}$. The map $j$ is a Thom-fiber inclusion: explicitly its source
before collapsing $EG$ is
$EG_+\wedge_G(G_+\wedge S^{-V})\simeq EG_+\wedge S^{-2}$,
and restriction to any point of $EG$ includes the corresponding
fiber in the Thom colimit. This also recovers the usual
circle-transfer sequence; see the first displayed fibration in
\cite[\S1]{EH5MadsenSchlichtkrull2000}.

\emph{The unit extends when $\eta_R=0$.}
Let $u:S^{-2}\to\Sigma^{-2}R$ be the shifted unit. Its obstruction to
extension over $j$ is the composite
\[
 \Sigma^\infty BG_+\xrightarrow{\delta}S^{-1}
      \xrightarrow{\Sigma u}\Sigma^{-1}R.
 \tag{D.2.14d}
\]
Since $\Sigma^\infty BG_+\simeq S\vee D$ and $F(D,R)=0$,
restriction to the basepoint gives an isomorphism
\[
 [\Sigma^\infty BG_+,\Sigma^{-1}R]
       \xrightarrow{\ \cong\ }[S,\Sigma^{-1}R].
\]
The restriction of \textup{(D.2.14d)} is the image under the unit of
an element of $\pi_1S\cong\mathbb Z/2$, generated by $\eta$.
It is therefore zero when $\eta_R=0$. Exactness of mapping out of
\textup{(D.2.14c)} supplies
\[
 t:\operatorname{Th}(-L)\longrightarrow\Sigma^{-2}R,
 \qquad t j=u.
 \tag{D.2.14e}
\]

\emph{The extension trivializes the full local system.}
Extend $t$ $R$-linearly to a map of left $R$-modules
$R\wedge\operatorname{Th}(-L)\to\Sigma^{-2}R$.
The Thom object is the colimit of the local system
$R\wedge S^{-L}$ of left $R$-modules over $BG$.
By the colimit adjunction, this map is the same as a natural
transformation
\[
 R\wedge S^{-L}\longrightarrow
          \operatorname{const}_{BG}(\Sigma^{-2}R).
 \tag{D.2.14f}
\]
Its value at the selected basepoint is the identity of
$\Sigma^{-2}R$, by \textup{(D.2.14e)} and the identification of $j$.
Every point of $BG$ can be joined to that basepoint. Naturality along
a path conjugates the corresponding fiber maps by equivalences.
Thus every component of \textup{(D.2.14f)} is an equivalence, and the
natural transformation itself is an equivalence of local systems.
This is the Thom-orientation criterion of
\cite[Definitions 2.20 and 2.22, Theorem 2.24]{EH5ABGHR2014}; it
retains the entire coherent sphere-bundle action.

For the stated $E_1$ generality, dualize pointwise using
$F_R(-,R)$ from left to right $R$-modules. Sphere lines are
dualizable, so \textup{(D.2.14f)} gives an equivalence of right
$R$-module local systems
\[
 S^L\wedge R\simeq
          \operatorname{const}_{BG}(\Sigma^2R).
\]
This step needs no commutative multiplication on $R$.
Taking colimits and then the augmentation $BG_+\to S$ produces a
map $\operatorname{Th}(L)\wedge R\to\Sigma^2R$ whose restriction
to the selected fiber is the identity. Precomposition with the unit
therefore gives
\[
 D\simeq\operatorname{Th}(L)\longrightarrow\Sigma^2R
\]
restricting to the unit on the bottom Thom sphere $S^2$.
But $F(D,R)=0$ makes this map null. The unit of $R$ is consequently
null, so unitality gives $\operatorname{id}_R=0$ and hence $R=0$.
This proves the first statement.

\emph{The exact local-unit criterion.}
Every localizing subcategory of spectra is a tensor ideal: for a
fixed object $A$ in the subcategory, the spectra $X$ with
$A\wedge X$ in that subcategory form a localizing subcategory
containing $S$, hence all spectra. Since $J_p$ is generated by one
object, it is the kernel of an accessible exact localization $L$.
Its equivalences are preserved by tensoring with arbitrary spectra.
The localized category therefore has its localized symmetric
monoidal structure, and its inclusion in spectra is lax symmetric
monoidal. Its unit gives an $E_\infty$-ring $R_0=L(S)$; this does
not assert that $L$ is smashing. These are precisely the monoidal
localization statements in
\cite[Example 2.2.1.7, Proposition 2.2.1.9]{FGHLurieHA}.

If $J_p$ is proper in $\operatorname{loc}(S/p)$, then
\[
 R_0/p\simeq L(S/p)\ne0.
\]
Let $C=\operatorname{fib}(S\to S[1/p])$. For any spectrum $E$,
the fiber sequence
\[
 F(S[1/p],E)\longrightarrow E\longrightarrow F(C,E)
\]
exhibits $F(C,E)$ as its derived $p$-completion: the first term is
$p$-invertible, and the last is right orthogonal to $S[1/p]$ because
$C\wedge S[1/p]=0$. This is orthogonality to all $p$-invertible
spectra, which are generated by $S[1/p]$.
In particular $F(C,E)/p\simeq E/p$.
Derived $p$-completion is a monoidal localization, since
$S/p$-equivalences are tensor-stable. Hence
\[
 R=\widehat{R_0}_p\simeq F(C,R_0)
\]
is an $E_\infty$-ring, with $R/p=R_0/p\ne0$.
It is still $J_p$-local, since
\[
 F(D\wedge S/p,F(C,R_0))
   \simeq F(D\wedge S/p\wedge C,R_0)=0.
\]
It follows that $p$ acts invertibly on $F(D,R)$.
This function spectrum is itself derived $p$-complete: if $X$ is
$p$-invertible, then
$F(X,F(D,R))\simeq F(X\wedge D,R)=0$.
Being both $p$-invertible and derived $p$-complete, it is zero.
Thus a proper inclusion produces the required nonzero ring $R$.

Conversely, if such a ring $R$ exists, then
$F(D\wedge S/p,R)=0$, so $R$ is $J_p$-local.
Equality $J_p=\operatorname{loc}(S/p)$ would make multiplication
by $p$ invertible on $R$. Its derived $p$-completeness would then
force $R=0$, a contradiction. This proves the second statement.
No multiplicative structure on the spectral cofiber $R_0/p$ has
been assumed.

\emph{Odd primes and the Pr\"ufer comparison.}
If $p$ is odd, multiplication by $2$ is invertible on every derived
$p$-complete spectrum. Hence $\eta_R=0$, and the first statement
excludes the ring in the second. Thus $J_p=\operatorname{loc}(S/p)$.

For completeness, write $P=\Sigma^\infty BC_{p^\infty}$.
The abelian-group extension
$0\to\mathbb Z\to\mathbb Z[1/p]\to C_{p^\infty}\to0$
gives a fiber sequence
\[
 K(\mathbb Z[1/p],1)\longrightarrow BC_{p^\infty}
       \longrightarrow K(\mathbb Z,2).
\]
The first space is the telescope of the degree-$p$ maps on $S^1$,
so it is mod-$p$ acyclic. The map $P\to D$ is therefore an
$H\mathbb F_p$-equivalence. Its cofiber $Q$ is connective.
The groups $H_*(Q;\mathbb Z)$ have multiplication by $p$ invertible,
by the mod-$p$ coefficient sequence. Consequently $Q/p$ has zero
integral homology, and stable Hurewicz for a bounded-below spectrum
gives $Q/p=0$. Thus $P/p\simeq D/p$.
This is also the Pr\"ufer-to-circle comparison used in
\cite[proof of Theorem 1.5]{EH5Levy2025}.

The reduced homology of $BC_{p^\infty}$ is $p$-primary, as follows
either from its filtered union of finite cyclic groups or from the
same extension. Hence $P[1/p]=0$, by stable Hurewicz.
The cofiber of $P\to P[1/p]$ is the filtered homotopy colimit of
the cofibers $P/p^r$ and is $\Sigma P$.
Each $P/p^r$ is a finite extension of copies of $P/p$; conversely
$P/p$ lies in $\operatorname{loc}(P)$.
Therefore
\[
 \operatorname{loc}(P)=\operatorname{loc}(P/p)
     =\operatorname{loc}(D/p)=J_p,
\]
which completes \textup{(D.2.14a)}.
\end{proof}

At $p=2$, the second statement is still an exact criterion, and
the first identifies the possible obstruction $\eta_R\ne0$.
D.2.15 removes this obstruction using the reduced circle transfer
and the endomorphism ring of $R\wedge C\eta$, proving the equality
at $2$ as well. The action of $\eta$ on $C\eta$ itself is nonzero;
the transfer is what makes its coefficient version vanish.
